% Week 10 实验周报 —— 编译方式：XeLaTeX
\documentclass[11pt]{ctexart}
\usepackage[a4paper, margin=2.5cm]{geometry}
\usepackage{booktabs}
\usepackage{hyperref}
\usepackage{enumitem}
\setlist{nosep}

\title{\textbf{Chooser Option 定价项目 —— 第 10 周实验报告}\\
\large Enhancement 第 2 周：S\&P 500 股票池、批量流水线、优化实验与面板检验}
\author{佟铸城（Zhucheng Tong）}
\date{2026 年 9 月 29 日}

\begin{document}
\maketitle

\section{本周目标}
把第 9 周的单票在线骨架扩到 S\&P 500 历史成分股；按导师"补充实验"的要求把优化做出来并给出优化前后倍数；兑现两条承诺（向量化批量定价本周接上、日更压到分钟级）；理论线在面板上复检 GARCH 锚并做锚定 ML 第一轮。

\section{完成工作}
\begin{enumerate}
  \item \textbf{股票池}：从 Wikipedia 成分表与历史变更表（公开来源）逆向回放出 2018 年至今的点时成员区间，666 只标的、任意日期 505--509 只；590 只有价格历史，76 只已退市标的 Yahoo 无历史，如实标注；
  \item \textbf{多源采集}：yfinance 批量 $\to$ 单票 $\to$ Alpha Vantage 三级回退，指数退避重试；本周 590 只全部由 yfinance 批量拿到，未触发回退；
  \item \textbf{批量流水线}：增量拉取按所需窗口分组、特征增量重建（尾部 760 行重算、回写 40 行，EWMA 从存量播种）、EWMA 与利率偏差落库、一次向量化全池定价与情景表；
  \item \textbf{前端}：新增"全市场概览"页签（578 只已定价标的、行业汇总）；
  \item \textbf{优化实验 E1--E5}：同一进程复测"前"路径，倍数 5.4 到 164；590 只日更实测 38 秒；
  \item \textbf{理论线}：GARCH 期限结构锚在 497 只标的上复检；锚定残差 ML 第一轮（面板、按日期净化切分、季度 walk-forward、测试集封存）；
  \item \textbf{自检}：十项审计全过（含第 9 周八项回归）；确定性重跑、报告数字回查、前端两种模式首屏、依赖扫描、打包泄露检查。
\end{enumerate}

\section{关键结果}
\begin{itemize}
  \item \textbf{优化}：特征重建 9.5$\times$，报价状态加载 6.8$\times$，全池报价 5.4$\times$，全池情景 164$\times$，日更流水线相对第 9 周单票串行外推 70$\times$（2654 s $\to$ 38 s）。"分钟级"承诺兑现且为秒级。剩余大头是四个串行的指数与利率请求（12.8 s）和逐票特征写入（14.7 s）；全池报价的 412 ms 几乎全在 \texttt{DISTINCT ON} 扫 124 万行，索引是明确的下一步。
  \item \textbf{GARCH 锚在面板上没有站住}：第 9 周 JPM 上 0.804 的综合分，在 497 只标的上变成胜率 54.5\%、中位 0.952，分布双峰（60 只长尾 $>2.5$，持续性接近 1 且全部正偏差）。GARCH 1 日锚反而更稳（胜率 72\%、池化综合分 0.812、偏差近零）。决定：\textbf{不换锚}，1 日锚列为第 11 周第二候选，期限结构锚只在持续性远离 1 时有意义。
  \item \textbf{锚定 ML 第一轮全部劣于 EWMA}：验证集（2023-06 至 2024-03，95{,}345 行、494 只）Ridge 综合分 7.16、GBDT 2.23，方向一致地高估（Ridge 偏差 $+27$ 美元）。特征截尾只改善到 6.12，说明主因是时间上的分布漂移而非横截面离群值——把样本从 1 只扩到 494 只并没有消除第 5 周的机制。第二轮方向：按日期横截面去均值、固定效应、GARCH 1 日锚；测试集仍封存。
\end{itemize}

\section{实验记录与发现}
\begin{itemize}
  \item 多标的暴露了单票看不到的三类脏数据：无分红标的的股息同比 $0/0$、首次分红的 $+\infty$、坏价格行的 Parkinson 无穷大，以及 AIV 因清算分红出现 228\% 股息率。处理为：股息同比在上年为零时定义为 0，Parkinson 仅对正价格计算，定价用股息率上限 25\%（原值保留），非有限输入的标的从批次剔除。JPM 文件特征逐位不变。
  \item 审计 6 初版用严格不等式，在 AIV 上因 call 腿为零、chooser 与 put 浮点相等而失败；这不是套利，改为带容差的 $\le$。
  \item \textbf{第 9 周审计 1 与 5 本周失败}，根因是 Yahoo 在每次分红后重算整段复权收盘价，本地快照与在线库是两个数据版本。改为同版本比对（用第 2 周文件管道在库内价格上重算，与库内特征逐位一致），历史报价行按"当日审计留痕"对待、只断言最新一行可复现。第 4 周存档审计失效的原因也由此完全解释。
  \item 一个曾被我的取最小值写法掩盖的问题：批量拉取若用整批的最早"上次日期"作窗口，一个陈旧标的会让全部 590 只重拉一个多月，覆盖掉库内旧版本。改为按窗口分组。
  \item 工具概览页初版对非 JPM 标的展示的 ML 参考价是分布外外推（43 只触到裁剪边界）。经导师确认（9 月 23 日），服务层加入验证范围门控：ML 参考价只对 JPM 计算和展示，其余标的字段为空、只展示主定价；审计 4 断言这一边界。替换为面板模型的前提是第 11 周第二轮在验证集上赢过基线。
\end{itemize}

\section{下周计划（第 11 周）}
工程线：Russell 3000 扩池；特征表索引或最新行视图；指数与利率请求并发。理论线：锚定 ML 第二轮（横截面去均值、固定效应、GARCH 1 日锚），测试集开封一次并与 EWMA / GARCH 1 日双基线比较；SHAP；OptionMetrics 隐含波动率初检（仅本机）。

\end{document}
